\section{Design and Architecture}\label{sec:design}

\subsection{Layers}

The program is one Maven module whose code sits under the base package \texttt{io.mcs.taskmanager}, divided by
responsibility into five packages (Fig.~\ref{fig:layers}), the last of them with a sub-package of its own,
\texttt{ui.kit}. Each has a \texttt{package-info.java} that says what it
holds:
\begin{itemize}
  \item \texttt{entities}: the task and its status, mapped by Hibernate onto the database table;
  \item \texttt{repositories}: Spring Data JPA interfaces, which Spring implements when the program starts;
  \item \texttt{services}: the operations the program offers on tasks, and every rule a change must pass;
  \item \texttt{ui}: the JavaFX window, made of FXML views, their controllers and a stylesheet;
  \item \texttt{shared}: what the other packages need but none of them owns, such as where the data folder is.
\end{itemize}
One rule holds between the layers: the screens call services, never repositories. The classes are few:
\begin{itemize}
  \item \texttt{TaskManagerApplication}, in the base package: the Spring Boot application class, whose
        \texttt{main} method starts the program;
  \item \texttt{entities.Task} and \texttt{entities.TaskStatus}: the task and its status (Section~\ref{sec:dsa});
  \item \texttt{repositories.TaskRepository}: the queries on the \texttt{tasks} table;
  \item \texttt{services.TaskService}: every operation on tasks; \texttt{services.TaskCounts}: the number of tasks
        in each status; \texttt{services.Refused}: the exception a refused change throws;
  \item \texttt{shared.AppDataFolder}: finds, and on the first start makes, the folder the database is kept in;
  \item \texttt{ui.DesktopApp}: the JavaFX application; \texttt{ui.MainWindowController}: the window, with
        \texttt{main-window.fxml} and the stylesheet \texttt{taskmanager.css}; \texttt{ui.TaskCell}: one row of the
        list; \texttt{ui.TaskDialog} and \texttt{ui.DeleteTaskDialog}: the task form and the delete confirmation,
        shown by \texttt{ui.DialogFrame}; \texttt{ui.TaskText}: the words the screen puts together, such as how a
        time is written; \texttt{ui.Dialog}, what every dialog has in common; \texttt{ui.StatusLook}, the icon
        and colour of each status; and \texttt{ui.kit.Ui} and \texttt{ui.kit.Icons}, small helpers that build
        the controls and draw the icons.
\end{itemize}

\begin{figure}[!t]
  \centering
  \begin{tikzpicture}[
      node distance=3.5mm and 4mm, >=Latex, font=\footnotesize,
      layer/.style={draw, rounded corners=2pt, minimum width=48mm, minimum height=8mm, align=center},
      side/.style={draw, dashed, rounded corners=2pt, text width=21mm, minimum height=8mm, align=center}]
    \node[layer] (ui) {\texttt{ui}: FXML views, controllers, CSS};
    \node[layer, below=of ui] (services) {\texttt{services}: operations and rules};
    \node[layer, below=of services] (repositories) {\texttt{repositories}: Spring Data JPA};
    \node[layer, below=of repositories] (db) {SQLite file \texttt{taskmanager.db}\\(schema by Flyway)};
    \node[side, right=of services, yshift=-5.75mm, minimum height=19.5mm] (entities)
      {\texttt{entities}:\\the task and its status};
    \node[side, right=of db] (shared) {\texttt{shared}:\\the data folder};
    \draw[->] (ui) -- node[right] {calls} (services);
    \draw[->] (services) -- node[right] {calls} (repositories);
    \draw[->] (repositories) -- node[right] {Hibernate} (db);
    \draw[->, dashed] (services.east) -- (services.east -| entities.west);
    \draw[->, dashed] (repositories.east) -- (repositories.east -| entities.west);
    \draw[->, dashed] (shared) -- (db);
    \draw[->, dashed] (ui.east) -| (entities.north);
  \end{tikzpicture}
  \caption{The layers of the program, in package \texttt{io.mcs.taskmanager}. Screens call services, never
           repositories; a dashed arrow means ``uses''.}\label{fig:layers}
\end{figure}

\subsection{Spring under JavaFX}

The program is a JavaFX application running on top of a Spring Boot application context with no web server. The
main class, \texttt{TaskManagerApplication}, is the Spring Boot application class, but it does not itself extend
JavaFX's \texttt{Application}: when the program is started from a plain class path, as Maven and the IDE start it, a
main class that extends \texttt{Application} stops with an error saying that the JavaFX runtime components are
missing. Its \texttt{main} method therefore hands over to a separate JavaFX application class, \texttt{ui.DesktopApp},
through \texttt{Application.launch}. JavaFX then calls three of its methods in turn:
\begin{itemize}
  \item \texttt{init()} starts Spring, before the window opens, with the data folder from \texttt{AppDataFolder}, so
        that the database has been found or made and migrated by Flyway by then;
  \item \texttt{start(Stage)} loads the window's FXML with Spring as the FXML loader's controller factory, so the
        window's controller is a Spring bean, then applies the stylesheet and shows the window;
  \item \texttt{stop()} closes Spring once the window has closed, and with it the database.
\end{itemize}
Because the controller is a Spring bean, it is handed the task service through its constructor, like any other bean,
and never creates it itself.

\subsection{Persistence}

Tasks are kept in an SQLite database~\cite{sqlite}, a single file that needs no server and no installation:
\path{%LOCALAPPDATA%\TaskManager\taskmanager.db}. The folder belongs to the signed-in Windows user and can be written
without administrator rights, and it is local rather than roaming, so the database is not copied around with the
user's profile.

The schema is created by Flyway~\cite{flyway} from versioned SQL scripts in
\path{src/main/resources/db/migration}, which run every time the program starts and apply any script the database
has not yet seen. Hibernate~\cite{hibernate} maps the task entity onto the table the scripts made, and only checks
that the two match (\texttt{spring.jpa.hibernate.ddl-auto=validate}); it never changes the schema itself. Spring
Data JPA~\cite{spring-boot} generates the repository's implementation from its interface.

Two copies of the program may have the same file open, and the connection is set up for it (Listing~\ref{lst:url}).
The database is opened in SQLite's write-ahead-log mode, which lets one copy read the file while the other writes.
Every transaction begins by taking the write lock (\texttt{transaction\_mode=IMMEDIATE}), so a copy that has read a
task saves its change before the other copy can change or delete that task, and a write that finds the lock taken
waits up to five seconds (\texttt{busy\_timeout=5000}) instead of failing at once. Each copy uses a single
connection, since SQLite takes one writer at a time anyway. This is what makes the two-copy scenario of
Section~\ref{sec:problem} safe to run.

\source[firstline=11, lastline=14, language={}]{src/main/resources/application.properties}
  {The connection, in \texttt{application.properties}}{lst:url}

\subsection{Rules and errors}

Every rule a change must pass is checked in a service: that the title is not blank once trimmed, that the title and
the description are within their length limits, and that the task still exists. A change the program will not make
is refused with one exception type, \texttt{Refused}, an unchecked exception whose message is written for the user,
such as ``A task needs a title, such as "Submit the OEL report".'' The service throws it inside its transaction, so
nothing of the refused change is saved; the screen catches it and shows the message as it is.

Keeping the rules in the services means that each rule lives in exactly one place, that the screens stay free of
rules, and that the JUnit tests exercise the same rules the user meets, without opening a window.

The entity's Bean Validation annotations stay as a last check behind the service, but they are not a way to refuse a
change: Hibernate checks them when it writes the task, which here is when the transaction commits, and a failure
there arrives as a \texttt{TransactionSystemException} with no message meant for the user. This is why the length limits are checked in
the service, with their own messages. Before they were, a title of 101 characters got past the service and failed
only at the commit (Section~\ref{sec:git} shows the change that fixed it).
An error the service does not expect, which is not a \texttt{Refused}, is caught by the screen, written to the log,
and shown as one plain sentence; how the user meets each kind of error is described in
Section~\ref{sec:functionality}.

\subsection{Reasons for the design}

\paragraph*{JavaFX rather than Swing} The brief allows either. JavaFX~\cite{openjfx} keeps the layout of each view in an FXML file
and its look in a CSS stylesheet, apart from the Java code, and its table and list controls are bound to observable
lists that redraw themselves when the data changes.

\paragraph*{Spring rather than objects wired by hand} With Spring, services and controllers are handed their
collaborators instead of constructing them, each service operation runs in a transaction, and the tests start the
same application context, with the same services and rules, on a database of their own.

\paragraph*{SQLite rather than a text file} A database gives transactions, so a change is saved completely or not at
all, and it lets two copies of the program share the data safely. SQLite gives this without a server.

\paragraph*{Flyway rather than letting Hibernate create the tables} The schema is versioned together with the code,
and every database, the user's and each test's, is built by the same scripts. Because Hibernate only validates, a
mismatch between the code and the schema stops the program when it starts instead of surfacing later as a failed
save.

\paragraph*{No code generation for accessors} Every constructor, getter and setter is written out by hand rather than
generated by a library such as Lombok, because the Javadoc tool reads the source: only members that are in the source
can be documented.
